\chapter{O Perfil do Consumidor do BNDES}

O perfil dos clientes do Banco Nacional de Desenvolvimento Econômico e Social (BNDES) é bastante diversificado. Diferentemente de instituições financeiras que atuam precipuamente com o público consumidor final, o BNDES atende diferentes tipos de pessoas físicas e jurídicas que buscam apoio financeiro para viabilizar investimentos produtivos e expandir suas operações. Entre os clientes elegíveis estão empresas privadas de diversos portes sediadas no país, associações, fundações, empresários individuais, pessoas físicas (como microempreendedores individuais, produtores rurais e transportadores autônomos de carga), além de órgãos e entidades da Administração Pública direta e indireta. Existem ainda hipóteses específicas de atendimento a empresas sediadas no exterior, desde que observadas as condições e requisitos normativos do Banco \cite{bndes2026cliente}.

As empresas são enquadradas pelo BNDES segundo a sua Receita Operacional Bruta (ROB) anual, dividindo-se nas categorias de micro, pequena, média e grande empresa. Essa estratificação baliza as condições financeiras, taxas e modalidades de apoio disponibilizadas pela instituição \cite{bndes2026cliente}.

Para identificar a distribuição do crédito concedido, foram avaliados os dados de desembolsos do BNDES consolidados em 2025 \cite{bndes_estatisticas_desembolsos_2026}. Durante o período, foram liberados aproximadamente R\$~169,7 bilhões para pessoas jurídicas de variados portes. Desse montante, as grandes empresas absorveram cerca de R\$~102,5 bilhões, o que representou 60,4\% dos desembolsos totais. As médias empresas responderam por R\$~38,1 bilhões (22,5\%), ao passo que as pequenas empresas obtiveram R\$~18,7 bilhões (11,0\%) e as microempresas R\$~10,3 bilhões (6,1\%) \cite{bndes_estatisticas_desembolsos_2026}.

\begin{table}[htbp]
  \centering
  \caption{Desembolsos do BNDES por porte da empresa em 2025}
  \label{tab:desembolsos_porte_2025}
  \begin{tabular}{lrr}
    \hline
    \textbf{Porte da empresa} & \textbf{Desembolsos em 2025} & \textbf{Participação} \\
    \hline
    Micro   & R\$ 10,3 bilhões  & 6,1\%  \\
    Pequena & R\$ 18,7 bilhões  & 11,0\% \\
    Média   & R\$ 38,1 bilhões  & 22,5\% \\
    Grande  & R\$ 102,5 bilhões & 60,4\% \\
    \hline
    \textbf{Total} & \textbf{R\$ 169,7 bilhões} & \textbf{100,0\%} \\
    \hline
  \end{tabular}
  \fonte{Elaborado pelos autores com base no Portal de Dados Abertos do \citeonline{bndes_estatisticas_desembolsos_2026}.}
\end{table}

Sob a perspectiva financeira do volume liberado, constata-se a preponderância das grandes organizações na destinação dos recursos em 2025. Contudo, essa concentração em valores monetários não implica necessariamente hegemonia no número absoluto de tomadores ou contratos celebrados, haja vista que as estatísticas do Portal de Dados Abertos desagregam os desembolsos por métricas como porte, setor econômico e localização geográfica \cite{bndes_estatisticas_desembolsos_2026}.

No tocante à distribuição regional, as operações em 2025 evidenciaram forte concentração geográfica nas regiões Sul e Sudeste: a Região Sudeste concentrou aproximadamente 40,0\% dos desembolsos, ao passo que a Região Sul recebeu 33,4\%. O Centro-Oeste alcançou 13,1\%, o Nordeste 9,3\% e a Região Norte respondeu por 4,2\% do montante \cite{bndes_estatisticas_desembolsos_2026}.

\begin{table}[htbp]
  \centering
  \caption{Distribuição dos desembolsos do BNDES por região em 2025}
  \label{tab:desembolsos_regiao_2025}
  \begin{tabular}{lr}
    \hline
    \textbf{Região} & \textbf{Participação aproximada} \\
    \hline
    Sudeste      & 40,0\% \\
    Sul          & 33,4\% \\
    Centro-Oeste & 13,1\% \\
    Nordeste     & 9,3\%  \\
    Norte        & 4,2\%  \\
    \hline
  \end{tabular}
  \fonte{Elaborado pelos autores com base no Portal de Dados Abertos do \citeonline{bndes_estatisticas_desembolsos_2026}.}
\end{table}

Conforme indicado na Tabela~\ref{tab:desembolsos_regiao_2025}, as regiões Sudeste e Sul responderam, conjuntamente, por 73,4\% dos desembolsos do BNDES no exercício. Estados como São Paulo, Rio Grande do Sul, Paraná, Santa Catarina e Minas Gerais figuraram entre os principais polos de alocação dos aportes, refletindo a densidade econômica e a estrutura industrial desses estados \cite{bndes_estatisticas_desembolsos_2026}.

No que se refere às demandas atendidas, o Banco estrutura linhas e programas operacionais orientados a múltiplas finalidades, como aquisição e produção de máquinas e equipamentos credenciados, projetos de implantação, modernização ou expansão da capacidade instalada, reforço de capital de giro associado, suporte a projetos de infraestrutura e iniciativas de pesquisa, desenvolvimento e inovação tecnológica \cite{bndes2026faq}. 

Em síntese, o perfil dos demandantes de recursos do BNDES abrange desde o microempreendedor até conglomerados e entes federativos, convergindo uma carteira ampla de clientes \cite{bndes2026cliente}. Todavia, a análise quantitativa dos desembolsos revela que o volume financeiro em 2025 concentrou-se majoritariamente no segmento de grandes empresas e no eixo geográfico Sul-Sudeste \cite{bndes_estatisticas_desembolsos_2026}.